% =====================================================================
\chapter{Πειραματική Αξιολόγηση}
\label{ch:evaluation}
% =====================================================================

Η αξιολόγηση απαντά σε τέσσερα ερωτήματα:

\begin{description}
  \item[ΕΕ1 -- Ανίχνευση.] Εντοπίζει το \Probolos{} τις τυπικές μορφές επιθέσεων
        \en{USB} χωρίς ψευδώς θετικά σε συνηθισμένες συσκευές;
  \item[ΕΕ2 -- Συμπεριφορά.] Πώς ταξινομεί το στάδιο~3 διαφορετικά πρότυπα
        πληκτρολόγησης, συμπεριλαμβανομένων όσων μιμούνται άνθρωπο;
  \item[ΕΕ3 -- Κόστος.] Ποιο είναι το υπολογιστικό κόστος της ανάλυσης;
  \item[ΕΕ4 -- Ορθότητα.] Πόσο αξιόπιστη είναι η υλοποίηση και ποιες τεχνικές
        επαλήθευσης αποδείχθηκαν αποτελεσματικές;
\end{description}

\section{Μεθοδολογία και περιβάλλον}
\label{sec:eval-method}

Η αξιολόγηση έγινε κυρίως με \emph{εξομοιωμένες} συσκευές, για τρεις λόγους: (α)
η επαναληψιμότητα — κάθε σενάριο εκτελείται πανομοιότυπα όσες φορές χρειαστεί·
(β) η ασφάλεια — καμία πραγματική κακόβουλη συσκευή δεν συνδέεται σε σταθμό
εργασίας· και (γ) η πληρότητα — μπορούν να κατασκευαστούν αλλοιωμένοι
περιγραφείς που δεν θα παρήγαγε κανένα εμπορικό εργαλείο.

\begin{table}[htbp]
\centering
\caption{Περιβάλλον αξιολόγησης.}
\label{tab:env}
\begin{tabular}{@{}ll@{}}
\toprule
\textbf{Στοιχείο} & \textbf{Τιμή}\\
\midrule
Έκδοση λογισμικού      & \Probolos{} 1.0.0b2\\
Γλώσσα                 & \en{Python} 3.10–3.13 (μετρήσεις χρόνου: 3.11, \en{x86\_64})\\
Διανομή ανάπτυξης      & \en{Manjaro Linux} (\en{KDE Plasma})\\
Εξομοίωση συσκευών     & \en{dummy\_hcd} + \en{raw\_gadget} (πυρήνας \en{Linux})\\
Έλεγχοι ιδιοτήτων      & \en{Hypothesis}~\cite{maciver2019hypothesis}\\
Στατική ανάλυση        & \en{ruff}, \en{mypy}, \en{CodeQL}\\
\bottomrule
\end{tabular}
\end{table}

\section{Το πλαίσιο εξομοίωσης (\en{testbed})}
\label{sec:eval-testbed}

Ο κατάλογος \code{testbed/} δημιουργεί συσκευές \en{USB} λογισμικού και τις
«συνδέει» στον τοπικό πυρήνα μέσω \en{dummy\_hcd} και \en{raw\_gadget}. Για
κάθε προκαθορισμένο σενάριο (\en{preset}) υπάρχει και λειτουργία πρόβλεψης
(\code{--preview}) που αναλύει τους περιγραφείς τοπικά, χωρίς \en{root}, και
τυπώνει τι πρέπει να συμπεράνει το \Probolos. Η σύγκριση πρόβλεψης και
πραγματικής εκτέλεσης διαχωρίζει σφάλματα των κανόνων από σφάλματα της
απαρίθμησης.

\begin{lstlisting}[language=shell,caption={Επίδειξη κρίσιμου ευρήματος με εξομοιωμένη συσκευή \en{BadUSB}.},label=lst:demo]
sudo modprobe -a dummy_hcd raw_gadget
sudo python3 -m probolos --dry-run        # terminal 1: report only
sudo python3 -m testbed.spawn badusb      # terminal 2: plug in the fake
python3 -m testbed.spawn badusb --preview # no root: predicted verdict
\end{lstlisting}

\section{Σενάρια και αποτελέσματα ανίχνευσης (ΕΕ1)}
\label{sec:eval-scenarios}

Ο Πίνακας~\ref{tab:scenarios} συνοψίζει τα αποτελέσματα για τα έξι σενάρια του
πλαισίου εξομοίωσης. Σε όλα, η ετυμηγορία ταυτίζεται με την αναμενόμενη.

\begin{table}[htbp]
\centering
\small
\caption{Αποτελέσματα σταδίων 1–2 στα σενάρια εξομοίωσης.}
\label{tab:scenarios}
\begin{tabular}{@{}p{2.9cm}p{3.3cm}p{5.2cm}c@{}}
\toprule
\textbf{Σενάριο} & \textbf{Δηλώνει} & \textbf{Ευρήματα} & \textbf{Ετυμηγορία}\\
\midrule
\code{flashdrive}  & \en{Mass Storage} & — & \textcolor{okgreen}{\textbf{καθαρή}}\\
\code{keyboard}    & πληκτρολόγιο & \code{keyboard-at-high-speed} & \sevN\\
\code{overpowered} & πληκτρολόγιο, 800\,\en{mA} & \code{power-exceeds-bus-limit}, \code{keyboard-at-high-speed} & \sevW\\
\code{badusb}      & \en{Mass Storage} + πληκτρολόγιο & \code{storage-with-keyboard}, \code{self-contradictory-identity}, \code{multiple-distinct-functions}, \code{keyboard-at-high-speed} & \sevC\\
\code{drift-innocent} & \en{Mass Storage} & — (εγκρίνεται) & \textcolor{okgreen}{\textbf{καθαρή}}\\
\code{drift-weaponized} & ίδια ταυτότητα + πληκτρολόγιο & ό,τι το \code{badusb} \emph{και} \code{descriptor-drift} & \sevC\\
\bottomrule
\end{tabular}
\end{table}

\subsection{Παράδειγμα 1: «Στικάκι» που πληκτρολογεί}
\label{sec:example-badusb}

Το σενάριο \code{badusb} εμφανίζεται ως \en{«SanDisk Cruzer Blade»}, με
\en{VID:PID} \code{0781:5567} — αναγνωριστικά ενός πραγματικού, δημοφιλούς
στικακιού — αλλά δηλώνει δύο διεπαφές: μαζική αποθήκευση (\code{08/06/50}) και
πληκτρολόγιο εκκίνησης (\code{03/01/01}). Το Απόσπασμα~\ref{lst:report-badusb}
δείχνει την αναφορά που παράγει το \Probolos. Παρατηρούμε ότι:

\begin{itemize}
  \item η αναφορά ξεκινά με τις \emph{δηλωμένες λειτουργίες} (\en{Mass Storage}
        \textbf{+} \en{KEYBOARD}) και όχι με το όνομα της συσκευής·
  \item οι συμβολοσειρές της συσκευής εμφανίζονται σε εισαγωγικά, ως
        «μαρτυρία» της συσκευής και όχι ως κρίση του εργαλείου·
  \item τα ευρήματα ταξινομούνται κατά σοβαρότητα και καθένα συνοδεύεται από
        επεξήγηση σε απλή γλώσσα.
\end{itemize}

\begin{lstlisting}[caption={Αναφορά τερματικού για το σενάριο \code{badusb} (έξοδος της μονάδας \code{report.py}, χωρίς χρώματα).},label=lst:report-badusb]
NEW USB DEVICE · blocked, awaiting decision
──────────────────────────────────────────────────────────────
  Mass Storage (SCSI)
  + KEYBOARD
  "SanDisk" "Cruzer Blade"
  serial "TB-STICK-1"
  0781:5567 · bus 1 port 1-4 · 480 Mbps · 200 mA bus-powered

  ✖ CRITICAL · Storage device that can also type
      This device presents itself as both a disk and a
      keyboard. The storage half provides the alibi while the
      keyboard half can enter commands on your behalf.
      Legitimate flash drives never need a keyboard interface.
      This is the classic BadUSB layout.

  ⚠ WARNING · Calls itself storage, presents an input
  interface
      The device names itself 'SanDisk Cruzer Blade', which
      describes a storage product, yet it declares a keyboard
      interface. ...

  • NOTICE · Device combines several different functions
  • NOTICE · Keyboard on a high-speed connection
──────────────────────────────────────────────────────────────
  checked: identity + consistency. Stages 3 & 4 follow.
\end{lstlisting}

Επειδή υπάρχει εύρημα \sevC, η αποδοχή απαιτεί πληκτρολόγηση της λέξης
\code{authorize} στο τερματικό ή αναμονή 10\,\en{s} στον διάλογο, και η επιλογή
«πάντα» δεν προσφέρεται.

\subsection{Παράδειγμα 2: Νόμιμο στικάκι}

Η ίδια ταυτότητα χωρίς τη διεπαφή πληκτρολογίου (σενάριο \code{flashdrive}) δεν
παράγει κανένα εύρημα: «\code{✓ no inconsistencies in what it claims}». Στη
λειτουργία πράκτορα, μια τέτοια συσκευή χρειάζεται \emph{έναν} διάλογο. Στη
συνέχεια, το στάδιο~4 διαβάζει τον πίνακα διαμερισμάτων· ένα τυπικό στικάκι με
ένα διαμέρισμα \en{FAT32} ή \en{exFAT} από τον τομέα 2048 δεν παράγει ευρήματα.

\subsection{Παράδειγμα 3: Μετάλλαξη περιγραφέων}
\label{sec:example-drift}

Το σενάριο ελέγχει την ανίχνευση συσκευής που εγκρίθηκε και στη συνέχεια
«οπλίστηκε». Πρώτα συνδέεται η \code{drift-innocent} (μόνο αποθήκευση, σειριακός
\code{TB-DRIFT-1}) και ο χρήστης την εγκρίνει. Έπειτα συνδέεται η
\code{drift-weaponized}: ίδιος κατασκευαστής, προϊόν και σειριακός αριθμός, αλλά
με πρόσθετη διεπαφή πληκτρολογίου. Το κανονικοποιημένο αποτύπωμα $F(D)$
διαφέρει, άρα ενεργοποιείται ο \code{descriptor-drift} (\sevC) — \emph{επιπλέον}
των κανόνων συνέπειας. Ακόμη κι αν ένας επιτιθέμενος κατάφερνε να αποφύγει τους
κανόνες του σταδίου~2, η αλλαγή αυτή καθαυτή θα αρκούσε.

\begin{lstlisting}[language=shell,caption={Εκτέλεση του σεναρίου μετάλλαξης.},label=lst:drift]
sudo python3 -m testbed.spawn drift-innocent     # approve it
sudo python3 -m testbed.spawn drift-weaponized   # same serial, now a keyboard
sudo python3 -m probolos --history -v            # the ledger shows both
\end{lstlisting}

\subsection{Παράδειγμα 4: Εχθρικές συμβολοσειρές}

Μια συσκευή \en{HID} που ονομάζεται \code{"Logitech\textbackslash x1b[2J USB Keyboard"}
θα καθάριζε την οθόνη του τερματικού αν το όνομά της τυπωνόταν αυτούσιο, ενώ
ένας χαρακτήρας \code{U+202E} (\en{RIGHT-TO-LEFT OVERRIDE}) θα αντέστρεφε την
εμφάνιση του κειμένου που ακολουθεί~\cite{boucher2023trojansource}. Το
\Probolos{} εμφανίζει τέτοιους χαρακτήρες με διαφυγή (\code{\textbackslash x1b},
\code{\textbackslash u202e}) και, επειδή η συσκευή μπορεί να πληκτρολογεί,
ενεργοποιεί τον \code{crafted-strings-hid} (\sevC). Ο έλεγχος δεν αφορά μόνο
το τερματικό: ο καθαρισμός γίνεται στο σημείο εισόδου, ώστε και ένα μεταγενέστερο
\code{cat probolos.jsonl} να είναι ασφαλές.

\subsection{Παράδειγμα 5: Κάρτα μνήμης με κρυφό διαμέρισμα}

Με \code{--watch-media}, όταν εισάγεται κάρτα σε ήδη εγκεκριμένο αναγνώστη, το
\Probolos{} διαβάζει τη διάταξή της. Μια κάρτα με διαμέρισμα τύπου \code{0xEF}
(\en{EFI System}) ή με «κρυφό» τύπο \en{MBR} (π.χ. \code{0x1C}) παράγει
\code{media-efi-system-partition} ή \code{media-hidden-partition} (\sevC). Με
\code{--media-policy deauthorize}, ολόκληρος ο αναγνώστης απενεργοποιείται.

\section{Ανάλυση συμπεριφοράς πληκτρολόγησης (ΕΕ2)}
\label{sec:eval-behaviour}

Για να μελετηθεί η συμπεριφορά των κανόνων του σταδίου~3 ανεξάρτητα από το
υλικό, οι πραγματικές συναρτήσεις του \Probolos{} (\code{rules.behaviour\_findings})
τροφοδοτήθηκαν με συνθετικές παρατηρήσεις: ακολουθίες $n$ πατημάτων με
διαστήματα από κανονική κατανομή $\mathcal{N}(\mu, \sigma)$ και χρόνο πρώτου
πατήματος $t_0$. Κάθε σενάριο εκτελέστηκε 1000 φορές (σπόρος 42). Τα
αποτελέσματα φαίνονται στον Πίνακα~\ref{tab:behaviour}.

\begin{table}[htbp]
\centering
\small
\caption[Ευρήματα σταδίου 3 σε συνθετικά πρότυπα πληκτρολόγησης]{Ευρήματα σταδίου 3 σε συνθετικά πρότυπα πληκτρολόγησης (1000 επαναλήψεις
ανά σενάριο). \code{MGK}: \code{machine-generated-keystrokes},
\code{UT}: \code{unprompted-typing}, \code{UK}: \code{unexpected-keystrokes},
\code{IA}: \code{immediate-activity}.}
\label{tab:behaviour}
\begin{tabular}{@{}p{4.1cm}cccccccc@{}}
\toprule
\textbf{Πρότυπο} & $n$ & $\mu$ (\en{ms}) & $\sigma$ (\en{ms}) & $t_0$ (\en{s}) & \code{MGK} & \code{UT} & \code{UK} & \code{IA}\\
\midrule
\en{Ducky}, σταθερός ρυθμός   & 40 & 10  & 0{,}5 & 0{,}05 & 1000 & 0    & 0   & 1000\\
\en{Ducky} με τυχαιότητα       & 40 & 10  & 4     & 0{,}05 & 0    & 1000 & 0   & 1000\\
Ωφέλιμο φορτίο 30\,\en{ms}     & 20 & 30  & 9     & 0{,}05 & 46   & 954  & 0   & 1000\\
Μίμηση τύπου \en{Malboard}     & 40 & 150 & 60    & 0{,}30 & 0    & 1000 & 0   & 1000\\
Ανθρώπινος ρυθμός              & 40 & 180 & 80    & 0{,}40 & 0    & 1000 & 0   & 1000\\
Τυχαίο άγγιγμα (3 πλήκτρα)     & 3  & 120 & 50    & 1{,}20 & 1    & 0    & 999 & 0\\
\bottomrule
\end{tabular}
\end{table}

Τα αποτελέσματα επιβεβαιώνουν τη σχεδιαστική επιλογή της §\ref{sec:design-stage3}:

\begin{itemize}
  \item Ο κανόνας ρυθμού (\code{MGK}) πιάνει μόνο το «αφελές» ωφέλιμο φορτίο με
        σταθερό ρυθμό. Με μικρή τεχνητή τυχαιότητα ($\sigma=4$\,\en{ms}, $CV\approx0{,}4$)
        ή με μίμηση ανθρώπου, παρακάμπτεται πλήρως.
  \item Ωστόσο, \emph{κάθε} σενάριο με $n\ge5$ καταλήγει σε εύρημα \sevC{} μέσω
        του \code{UT}, ανεξάρτητα από τον ρυθμό. Η άμυνα δεν στηρίζεται στη
        διάκριση ανθρώπου--μηχανής, αλλά στο ότι κανείς δεν πρέπει να
        πληκτρολογεί.
  \item Το τυχαίο άγγιγμα τριών πλήκτρων αναφέρεται ως \sevW{} (\code{UK}) και
        όχι ως επίθεση, μειώνοντας τα ψευδώς θετικά· σε 1 από τις 1000
        επαναλήψεις τα τυχαία διαστήματα έτυχε να είναι αρκετά κανονικά ώστε να
        δώσουν \code{MGK}.
\end{itemize}

\begin{warnbox}[Περιορισμός]
Μια συσκευή που \emph{δεν} πληκτρολογεί κατά τη διάρκεια της καραντίνας (π.χ.
περιμένει 10\,\en{s}, ή ενεργοποιείται μόνο όταν ανιχνεύσει αδράνεια του
χρήστη) δεν εντοπίζεται από το στάδιο~3. Αυτό είναι εγγενές όριο κάθε
παρατήρησης περιορισμένου χρόνου και συζητείται στο Κεφάλαιο~\ref{ch:discussion}.
\end{warnbox}

\section{Υπολογιστικό κόστος (ΕΕ3)}
\label{sec:eval-cost}

Ο Πίνακας~\ref{tab:perf} δείχνει τον χρόνο ανάλυσης περιγραφέων και αξιολόγησης
κανόνων των σταδίων 1–2 ανά σενάριο (5000 επαναλήψεις, \en{Python} 3.11,
\en{x86\_64}). Ο διάμεσος χρόνος είναι 18–32\,\textmu\en{s} και το 99ο
εκατοστημόριο κάτω από 150\,\textmu\en{s}: η ανάλυση είναι αμελητέα σε σχέση με
τον χρόνο απαρίθμησης μιας συσκευής (εκατοντάδες \en{ms}) και πολύ περισσότερο
σε σχέση με τον χρόνο αντίδρασης του ανθρώπου.

\begin{table}[htbp]
\centering
\caption{Χρόνος σταδίων 1–2 (ανάλυση περιγραφέων + κανόνες).}
\label{tab:perf}
\begin{tabular}{@{}lrrc@{}}
\toprule
\textbf{Σενάριο} & \textbf{Διάμεσος (\textmu\en{s})} & \textbf{99\% (\textmu\en{s})} & \textbf{Ευρήματα}\\
\midrule
\code{flashdrive}       & 17{,}7 & 104{,}2 & 0\\
\code{keyboard}         & 22{,}2 & 119{,}0 & 1\\
\code{overpowered}      & 24{,}0 & 126{,}2 & 2\\
\code{badusb}           & 31{,}4 & 126{,}7 & 4\\
\code{drift-innocent}   & 18{,}5 & 116{,}0 & 0\\
\code{drift-weaponized} & 31{,}4 & 147{,}1 & 4\\
\bottomrule
\end{tabular}
\end{table}

Το κυρίαρχο κόστος βρίσκεται στα στάδια 3 και 4, και είναι \emph{σκόπιμο}: η
καραντίνα διαρκεί \code{--observe} δευτερόλεπτα (προεπιλογή 3), ενώ η ανάγνωση
αποθήκευσης εξαρτάται από το πότε η συσκευή εκθέτει τον κόμβο μπλοκ (σε δοκιμές,
1–3\,\en{s}). Ο χρήστης μπορεί να τα απενεργοποιήσει
(\code{--observe 0 --no-storage-scan}) για αυστηρότερη απομόνωση και μηδενικό
χρόνο αναμονής.

\paragraph{Κενό έκθεσης.} Σε πραγματικό υλικό, ο χρόνος από την εξουσιοδότηση
έως την πρώτη επιτυχή δέσμευση \code{EVIOCGRAB} μετρήθηκε στα
41–85\,\en{ms}. Ο χρόνος αυτός \emph{δεν} είναι ο αριθμός πλήκτρων που
διέφυγαν· ο ακριβής αριθμός απαιτεί το πείραμα διαρροής του
\code{testbed/hidexp} (§\ref{sec:eval-pending}).

\section{Επαλήθευση της υλοποίησης (ΕΕ4)}
\label{sec:eval-verification}

\subsection{Αυτοματοποιημένοι έλεγχοι}

Η σουίτα περιλαμβάνει 1413 ελέγχους μονάδας και ενσωμάτωσης (6 παραλείπονται
όταν δεν υπάρχουν \en{AF\_UNIX sockets}), που εκτελούνται σε περίπου 39\,\en{s}
χωρίς \en{root}, χωρίς υλικό και χωρίς \en{pyudev} (απαίτηση~\ref{req:testable}).
Το \en{sysfs} προσομοιώνεται με πραγματικούς καταλόγους και συμβολικούς
συνδέσμους σε προσωρινό σύστημα αρχείων (\code{FakeSysfs}). Οι έλεγχοι
παλινδρόμησης από ελέγχους κώδικα ονομάζονται με βάση το ελάττωμα που
«καρφιτσώνουν» (π.χ. \code{AnalyzerHasNoControllingTerminal}). Η κατανομή ανά
περιοχή φαίνεται στο Σχήμα~\ref{fig:tests}.

\begin{figure}[htbp]
\centering
\begin{tikzpicture}
\begin{axis}[
  xbar, width=12cm, height=9.5cm,
  xlabel={Πλήθος ελέγχων}, xmin=0,
  symbolic y coords={testbed,interrogate,service,properties,textsafe,sysfs,gate,descriptors,cli,trust,quarantine,rules,ledger,boundary,agent,daemon,privsep,storage},
  ytick=data, y tick label style={font=\scriptsize\ttfamily},
  nodes near coords, nodes near coords style={font=\scriptsize},
  bar width=6pt, enlarge y limits=0.04,
]
\addplot[fill=univblue!60, draw=univblue] coordinates {
 (6,testbed) (10,interrogate) (11,service) (30,properties) (34,textsafe)
 (49,sysfs) (50,gate) (51,descriptors) (56,cli) (71,trust) (77,quarantine)
 (85,rules) (115,ledger) (129,boundary) (160,agent) (164,daemon)
 (166,privsep) (172,storage)};
\end{axis}
\end{tikzpicture}
\caption{Συναρτήσεις ελέγχου ανά αρχείο της σουίτας (\code{tests/test\_*.py}).
Οι μεγαλύτερες ομάδες καλύπτουν τον διαχωρισμό προνομίων, την αποθήκευση και
τον δαίμονα.}
\label{fig:tests}
\end{figure}

\subsection{Έλεγχοι ιδιοτήτων}

Ο έλεγχος βασισμένος σε ιδιότητες (\en{property-based testing})~\cite{claessen2000quickcheck,maciver2019hypothesis}
δηλώνει \emph{αναλλοίωτες} ιδιότητες και αφήνει το πλαίσιο να παράγει εισόδους.
Παραδείγματα ιδιοτήτων του \Probolos:

\begin{itemize}
  \item για \emph{οποιαδήποτε} ακολουθία \en{bytes}, ο αναλυτής περιγραφέων δεν
        εγείρει εξαίρεση· είτε επιστρέφει δομή είτε σφάλμα ανάλυσης·
  \item για \emph{οποιαδήποτε} συμβολοσειρά, η έξοδος του \code{textsafe} δεν
        περιέχει χαρακτήρες ελέγχου, \en{bidi} ή αόρατους·
  \item η αναφορά για οποιαδήποτε συσκευή δεν περιέχει ακολουθίες διαφυγής.
\end{itemize}

Η «βαθιά» εκτέλεση (29 έλεγχοι ιδιοτήτων $\times$ 7500 παραδείγματα) διαρκεί
7–12 λεπτά και ήταν καθαρή σε όλες τις εκδόσεις από την 0.11.0 έως την 1.0.0b2.
Ήδη την πρώτη ημέρα εντόπισε ένα ελάττωμα που είχε διαφύγει κάθε άλλου ελέγχου:
η περικοπή συμβολοσειράς σε όριο 1 χαρακτήρα παρέλειπε τη διαφυγή ενός
χαρακτήρα ελέγχου.

\subsection{Ασαφής έλεγχος (\en{fuzzing})}

Ο κατάλογος \code{fuzzing/} περιέχει \en{fuzzer} καθοδηγούμενο από κάλυψη
(μόνο με τη βιβλιοθήκη προτύπων) για ό,τι αναλύει το \Probolos{} πριν από μια
απόφαση: περιγραφείς, αναφορές \en{HID}, μέσα αποθήκευσης, το πρωτόκολλο της
πύλης, την υποδοχή του πράκτορα, τα αρχεία κατάστασης και ολόκληρη τη διαδρομή
ανάλυσης. Ενώ οι έλεγχοι ιδιοτήτων δοκιμάζουν περίπου 150 εισόδους ανά
ισχυρισμό, ο \en{fuzzer} εκτελεί εκατομμύρια, κρατά όσες φτάνουν σε νέο κώδικα
και αναφέρει και \emph{κολλήματα} και \emph{εκρήξεις μνήμης}, όχι μόνο
εξαιρέσεις. Εντόπισε ένα ελάττωμα στον έλεγχο δομής \en{UDF}, το οποίο οδήγησε
σε στοχευμένη επανεξέταση όλων των αναλυτών μη έμπιστης εισόδου.

\subsection{Κάλυψη κώδικα}

Ο Πίνακας~\ref{tab:coverage} δείχνει την κάλυψη διακλαδώσεων της έμπιστης
βάσης. Μετά την πρώτη φάση σκλήρυνσης, η κάλυψη του ορίου προνομίων ανέβηκε από
80\% σε 95\%, και το κατώφλι στη συνεχή ενσωμάτωση (\en{CI}) αυξήθηκε αντίστοιχα.

\begin{table}[htbp]
\centering
\caption{Κάλυψη διακλαδώσεων της έμπιστης βάσης (ως \en{root}).}
\label{tab:coverage}
\begin{tabular}{@{}lrr@{}}
\toprule
\textbf{Μονάδα} & \textbf{0.11.0} & \textbf{Μετά τη σκλήρυνση}\\
\midrule
\code{protocol.py}     & 97\% & 98\%\\
\code{gate\_server.py} & 81\% & 95\%\\
\code{gate\_client.py} & 71\% & 94\%\\
\code{privsep.py}      & 67\% & 98\%\\
\code{securefs.py}     & 87\% & 87\%\\
\code{atomicio.py}     & 85\% & 85\%\\
\midrule
\textbf{Όριο προνομίων (σύνολο)} & \textbf{80\%} & \textbf{95\%}\\
Σύνολο έργου          & 81\% & 85\%\\
\bottomrule
\end{tabular}
\end{table}

\section{Ανάλυση του αρχείου ελαττωμάτων}
\label{sec:eval-qa}

Από το πάγωμα λειτουργιών (έκδοση 0.11.0) και μετά, κάθε ελάττωμα καταγράφηκε
στο \code{docs/QA-LOG.md} με ημερομηνία, περιγραφή, μέθοδο εντοπισμού και
προτεραιότητα. Η κλίμακα προτεραιότητας είναι:

\begin{description}
  \item[P0] αποτυχία-ανοιχτά (αποδοχή χωρίς ανθρώπινη απάντηση), πύλη ανοιχτή ή
        θύρες νεκρές μετά την έξοδο, κλιμάκωση προνομίων, αδυναμία εκκίνησης
        της υπηρεσίας·
  \item[P1] λάθος ετυμηγορία σε κοινό υλικό, χαμένη απόφαση, φθορά αποθηκών,
        αδυναμία ασφαλείας χωρίς γνωστή εκμετάλλευση·
  \item[P2] σπάνιες οριακές περιπτώσεις, παραπλανητικά μηνύματα, σφάλματα
        τεκμηρίωσης·
  \item[P3] αισθητικά ζητήματα.
\end{description}

Καταγράφηκαν συνολικά \textbf{106 ελαττώματα}: 4 P0, 15 P1, 47 P2 και 40 P3
(Σχήμα~\ref{fig:qa}). Επειδή το σύνολο λειτουργιών ήταν «παγωμένο», όλα αφορούν
το \emph{ίδιο} σύστημα, κάτι που επιτρέπει τη σύγκριση των μεθόδων εντοπισμού.

\begin{figure}[htbp]
\centering
\begin{tikzpicture}
\begin{axis}[
  ybar stacked, width=11.5cm, height=7cm,
  ylabel={Ελαττώματα}, ymin=0,
  symbolic x coords={Έλεγχος κώδικα, Καθημερινή χρήση, Ιδιότητες, Fuzzing, CI},
  xtick=data, x tick label style={font=\small, align=center, text width=2.2cm},
  xticklabels={Έλεγχος\\κώδικα, Καθημερινή\\χρήση, Έλεγχοι\\ιδιοτήτων, \en{Fuzzing}, \en{CI}},
  legend style={at={(0.98,0.97)}, anchor=north east, font=\small},
  bar width=22pt, enlarge x limits=0.14,
]
\addplot[fill=critical, draw=critical!70!black] coordinates {(Έλεγχος κώδικα,4) (Καθημερινή χρήση,0) (Ιδιότητες,0) (Fuzzing,0) (CI,0)};
\addplot[fill=warning, draw=warning!70!black] coordinates {(Έλεγχος κώδικα,15) (Καθημερινή χρήση,0) (Ιδιότητες,0) (Fuzzing,0) (CI,0)};
\addplot[fill=notice!70, draw=notice] coordinates {(Έλεγχος κώδικα,41) (Καθημερινή χρήση,4) (Ιδιότητες,1) (Fuzzing,0) (CI,1)};
\addplot[fill=gray!40, draw=gray] coordinates {(Έλεγχος κώδικα,36) (Καθημερινή χρήση,3) (Ιδιότητες,0) (Fuzzing,1) (CI,0)};
\legend{P0, P1, P2, P3}
\end{axis}
\end{tikzpicture}
\caption{Ελαττώματα ανά μέθοδο εντοπισμού και προτεραιότητα (106 συνολικά).}
\label{fig:qa}
\end{figure}

\paragraph{Παρατηρήσεις.}
\begin{enumerate}
  \item \textbf{Ο συστηματικός έλεγχος κώδικα κυριαρχεί} (96 από 106), και
        εντόπισε \emph{όλα} τα ελαττώματα P0 και P1. Οι έλεγχοι οργανώθηκαν
        ανά «γωνία» (όριο προνομίων· χειρισμός συσκευών και αναλυτές· ροή
        αποφάσεων· συσκευασία και τεκμηρίωση), και κάθε εύρημα επανελεγχόταν
        από ξεχωριστό, «σκεπτικιστή» αναθεωρητή πριν από τη διόρθωση.
  \item \textbf{Οι αυτοματοποιημένες τεχνικές βρήκαν λίγα αλλά μοναδικά
        ελαττώματα}: ό,τι εντόπισαν οι έλεγχοι ιδιοτήτων και το \en{fuzzing} δεν
        είχε εντοπιστεί από κανέναν έλεγχο κώδικα.
  \item \textbf{Η καθημερινή χρήση} ανέδειξε ζητήματα αλληλεπίδρασης με την
        επιφάνεια εργασίας (π.χ. αυτόματη προσάρτηση στο \en{KDE Plasma}) που
        καμία εξομοίωση δεν αναπαράγει.
\end{enumerate}

\paragraph{Τα ελαττώματα P0.} Δύο από αυτά είναι χαρακτηριστικά των κινδύνων
ενός εργαλείου που «κλείνει» το σύστημα:

\begin{itemize}
  \item Η ανάγνωση χαρακτηριστικών \en{sysfs} αποκωδικοποιούσε με την
        κωδικοποίηση του \en{locale}· σε ορισμένα συστήματα αυτό προκαλούσε
        \code{UnicodeDecodeError} κατά την εκκίνηση και η πύλη ξανάνοιγε καθώς
        η διεργασία τερμάτιζε — αθόρυβη απώλεια προστασίας.
  \item Η διεργασία ανάγνωσης του σταδίου~4 κληρονομούσε την οπλισμένη πύλη
        (§\ref{sec:impl-gate}): ένα \code{SIGTERM} μετά από κόλλημα την έκανε να
        ανοίξει όλους τους διανομείς, ενώ ο δαίμονας πίστευε την πύλη κλειστή.
        Μια εχθρική συσκευή αποθήκευσης θα μπορούσε να χρονομετρήσει το κόλλημα.
\end{itemize}

Και τα δύο διορθώθηκαν με ελέγχους παλινδρόμησης που αποτυγχάνουν στον κώδικα
πριν από τη διόρθωση και περνούν μετά.

\section{Μετρήσεις σε εξέλιξη}
\label{sec:eval-pending}

Για λόγους επιστημονικής ακεραιότητας, σημειώνεται ότι οι ακόλουθες μετρήσεις
\emph{δεν} έχουν ακόμη ολοκληρωθεί σε πραγματικό υλικό επίθεσης και
προγραμματίζονται για την ολοκλήρωση της εργασίας (Πίνακας~\ref{tab:pending}).
Οι κενές τιμές πρέπει να συμπληρωθούν από τον συγγραφέα πριν από την κατάθεση,
ή η ενότητα να μεταφερθεί στη μελλοντική εργασία.

\begin{table}[htbp]
\centering
\small
\caption{Προγραμματισμένες μετρήσεις σε πραγματικό υλικό (προς συμπλήρωση).}
\label{tab:pending}
\begin{tabular}{@{}p{5.6cm}p{4.6cm}p{3.0cm}@{}}
\toprule
\textbf{Μέτρηση} & \textbf{Μέθοδος} & \textbf{Αποτέλεσμα}\\
\midrule
Πλήκτρα που διαφεύγουν πριν από τη δέσμευση & \code{testbed/hidexp}, $\ge$30 επαναλήψεις & [—]\\
Ποσοστό ανίχνευσης \en{Rubber Ducky} / \en{O.MG} & πραγματικές συσκευές & [—]\\
Ψευδώς θετικά σε συνηθισμένες συσκευές & $\ge$30 εμπορικές συσκευές & [—]\\
Χρόνος σταδίων 3–4 από άκρο σε άκρο & χρονοσφραγίδες \en{JSON log} & [—]\\
Μήτρα υλικού (\en{KDE}/\en{GNOME}, \en{X11}/\en{Wayland}) & σενάρια εγκατάστασης/σύνδεσης & [—]\\
\bottomrule
\end{tabular}
\end{table}
